% T3 candidate (sonnet3) for fig:hysloop — main curve Omega=4RT unchanged in spirit (same
% eq:Veq evaluation as the existing figure), plus: (i) Maxwell potential line y=0 made explicit,
% (ii) numeric gap value from eq:dUhys, (iii) a curve FAMILY at the real Omega/RT ratios of the
% weakest (4->3) and strongest (2->1) graphite staging transitions from table tab:staging —
% not schematic, all three curves are eq:Veq evaluated at their own Omega (T3_calc.py).
\begin{tikzpicture}[x=9.2cm,y=1.3cm]
\draw[-{Latex[length=1.6mm]}] (0,-2.35) -- (0,2.35) node[left,font=\scriptsize] {$(sF/RT)(V_\eq-U_j)$};
\draw[-{Latex[length=1.6mm]}] (-0.02,0) -- (1.06,0) node[below,font=\scriptsize] {$\xi$};
\foreach \xx in {0.25,0.5,0.75} {\draw (\xx,0.03) -- (\xx,-0.03) node[below,font=\scriptsize] {\xx};}
\foreach \yy in {-2,-1,1,2} {\draw (0.006,\yy) -- (-0.006,\yy) node[left,font=\scriptsize] {\yy};}
% Maxwell potential (y=0, i.e. U_j itself in these normalized units)
\draw[densely dotted] (-0.02,0) -- (1.06,0);
\node[font=\tiny,align=left] at (0.09,0.14) {Maxwell potential $U_j$ ($y{=}0$)};
% ---- family: real graphite staging Omega/RT ratios (table tab:staging), thin comparison curves ----
\draw[gray,densely dashed] plot[smooth] coordinates {(0.0200,-1.5683) (0.0440,-0.8712) (0.0680,-0.5266) (0.0920,-0.3144) (0.1160,-0.1720) (0.1400,-0.0726) (0.1640,-0.0023) (0.1880,0.0473) (0.2120,0.0812) (0.2360,0.1032) (0.2600,0.1158) (0.2840,0.1209) (0.2916,0.1212) (0.3080,0.1199) (0.3320,0.1141) (0.3560,0.1043) (0.3800,0.0913) (0.4040,0.0759) (0.4280,0.0585) (0.4520,0.0398) (0.4760,0.0201) (0.5000,0.0000) (0.5240,-0.0201) (0.5480,-0.0398) (0.5720,-0.0585) (0.5960,-0.0759) (0.6200,-0.0913) (0.6440,-0.1043) (0.6680,-0.1141) (0.6920,-0.1199) (0.7084,-0.1212) (0.7160,-0.1209) (0.7400,-0.1158) (0.7640,-0.1032) (0.7880,-0.0812) (0.8120,-0.0473) (0.8360,0.0023) (0.8600,0.0726) (0.8840,0.1720) (0.9080,0.3144) (0.9320,0.5266) (0.9560,0.8712) (0.9800,1.5683)};
\draw[gray,densely dotted,thick] plot[smooth] coordinates {(0.0200,1.1426) (0.0440,1.7041) (0.0680,1.9131) (0.0920,1.9898) (0.1067,2.0001) (0.1160,1.9966) (0.1400,1.9605) (0.1640,1.8953) (0.1880,1.8093) (0.2120,1.7077) (0.2360,1.5942) (0.2600,1.4712) (0.2840,1.3408) (0.3080,1.2043) (0.3320,1.0629) (0.3560,0.9175) (0.3800,0.7690) (0.4040,0.6180) (0.4280,0.4651) (0.4520,0.3108) (0.4760,0.1556) (0.5000,0.0000) (0.5240,-0.1556) (0.5480,-0.3108) (0.5720,-0.4651) (0.5960,-0.6180) (0.6200,-0.7690) (0.6440,-0.9175) (0.6680,-1.0629) (0.6920,-1.2043) (0.7160,-1.3408) (0.7400,-1.4712) (0.7640,-1.5942) (0.7880,-1.7077) (0.8120,-1.8093) (0.8360,-1.8953) (0.8600,-1.9605) (0.8840,-1.9966) (0.8933,-2.0001) (0.9080,-1.9898) (0.9320,-1.9131) (0.9560,-1.7041) (0.9800,-1.1426)};
% ---- main curve: Omega=4RT (unchanged from existing figure) ----
\draw[thick] plot[smooth] coordinates {(0.0200,-0.0518) (0.0440,0.5694) (0.0680,0.8382) (0.0920,0.9745) (0.1160,1.0411) (0.1400,1.0647) (0.1464,1.0657) (0.1640,1.0592) (0.1880,1.0329) (0.2120,0.9911) (0.2360,0.9373) (0.2600,0.8740) (0.2840,0.8033) (0.3080,0.7265) (0.3320,0.6448) (0.3560,0.5592) (0.3800,0.4705) (0.4040,0.3792) (0.4280,0.2860) (0.4520,0.1914) (0.4760,0.0959) (0.5000,0.0000) (0.5240,-0.0959) (0.5480,-0.1914) (0.5720,-0.2860) (0.5960,-0.3792) (0.6200,-0.4705) (0.6440,-0.5592) (0.6680,-0.6448) (0.6920,-0.7265) (0.7160,-0.8033) (0.7400,-0.8740) (0.7640,-0.9373) (0.7880,-0.9911) (0.8120,-1.0329) (0.8360,-1.0592) (0.8536,-1.0657) (0.8600,-1.0647) (0.8840,-1.0411) (0.9080,-0.9745) (0.9320,-0.8382) (0.9560,-0.5694) (0.9800,0.0518)};
\fill (0.1464,1.0657) circle (0.5pt) node[above,font=\scriptsize] {$\xi_s^-$};
\fill (0.8536,-1.0657) circle (0.5pt) node[below,font=\scriptsize] {$\xi_s^+$};
% spinodal potential levels (horizontal guides) + gap
\draw[gray,densely dashed] (-0.02,1.0657) -- (1.06,1.0657);
\draw[gray,densely dashed] (-0.02,-1.0657) -- (1.06,-1.0657);
\draw[{Latex[length=1.4mm]}-{Latex[length=1.4mm]}] (1.0,-1.0657) -- (1.0,1.0657)
  node[pos=0.4,right,font=\scriptsize,align=left] {$\Delta U_j^\hys$\\[-1pt]{\tiny$=2.131\,RT/(sF)$}};
% direction arrows on overshoot paths (bold sub-paths of the SAME main curve, endpoints match exactly)
\draw[-{Latex[length=1.6mm]},very thick] plot[smooth] coordinates {(0.0200,-0.0518) (0.0440,0.5694) (0.0680,0.8382) (0.0920,0.9745) (0.1160,1.0411) (0.1400,1.0647) (0.1464,1.0657)};
\draw[-{Latex[length=1.4mm]},densely dashed] (0.1464,1.0) -- (0.1464,0.10);
\node[font=\scriptsize] at (0.32,1.55) {discharge overshoot};
\draw[-{Latex[length=1.6mm]},very thick] plot[smooth] coordinates {(0.8536,-1.0657) (0.8600,-1.0647) (0.8840,-1.0411) (0.9080,-0.9745) (0.9320,-0.8382) (0.9560,-0.5694) (0.9800,0.0518)};
\draw[-{Latex[length=1.4mm]},densely dashed] (0.8536,-1.0) -- (0.8536,-0.10);
\node[font=\scriptsize] at (0.68,-1.55) {charge overshoot};
\node[font=\scriptsize] at (0.20,-0.42) {$\gamma\!\to\!0$: plateau at $U_j$ ($y=0$)};
% in-plot curve-family labels at xi~0.66 where all three curves are well separated
\node[gray,font=\tiny] at (0.685,0.02) {$\Omega/RT{=}2.42$ ($4\!\to\!3$)};
\node[font=\tiny] at (0.685,-0.72) {$\Omega{=}4RT$ (main)};
\node[gray,font=\tiny] at (0.685,-1.20) {$\Omega/RT{=}5.24$ ($2\!\to\!1$)};
\end{tikzpicture}
\caption{비단조 평형 전위 $V_\eq(\xi)$(식~\eqref{eq:Veq}, 좌표는 식 그대로의 수치 평가 — 검증
T3\_calc.py)와 과주행. 실선(굵게) $=\Omega_j{=}4RT$(기존 그림과 동일 대응), 회색 곡선 둘은
표~\ref{tab:staging} 초기값 중 가장 작은 것($4\!\to\!3$, $\Omega/RT{=}2.42$, 파선)과 가장 큰
것($2\!\to\!1$, $\Omega/RT{=}5.24$, 점선)이다 — 같은 식을 다른 $\Omega$ 에 평가해 gap 이 $\Omega$
와 함께 커짐을 보인다(작을수록 문턱 $2RT$ 에 가까워 gap 이 좁고, 클수록 넓다). 점선 수평선 $y{=}0$ 이
Maxwell 전위 $U_j$; 화살표 달린 굵은 경로가 방전(좌상, $\xi_s^-$ 까지)$\cdot$충전(우하, $\xi_s^+$ 까지) 과주행
방향이다. 두 극값의 전위 차(수평 회색 점선 두 개 사이)가 spinodal 상한 gap $\Delta U_j^\hys$(식~\eqref{eq:dUhys},
$\Omega{=}4RT$ 일 때 $2.131\,RT/(sF)$). $\gamma\to0$ 가역 한계에서 두 경로가 $y=0$ plateau 로 합쳐진다.}
\label{fig:hysloop}
